\chapter{Programación de Restricciones}

% Satisfaccion de restricciones
Los problemas se componen por tres elementos:
\begin{itemize}
    \item Variables
    \item Dominio de las Variables
    \item Restricciones
\end{itemize}

Buscamos en profundidad porque buscamos una solucion, todas las soluciones van a tener la misma longitud. 
La anchura va creciendo en cada nivel por lo que tardaria mas a pesar de que daria todas las soluciones incluyendo la optima si existe, en principio no hay una optima.

% Euristicas
Primero lo que mas va a fallar:
- Una forma de optimizar los algoritmos es ir buscando las opciones que antes van a fallar para agilizar el arbol de busqueda un monton, por ejemplo buscar las que menos dominio tienen primero.
- Tambien podemos mirar la que tiene mayor dominio.
- La que participe en mas restricciones.
- A la hora de chequear las restricciones no lo mismo chequear la que este en restricciones mas pesadas que en lass que estan en mas livianas.
- ordenar las restricciones  de tal forma que se ordenen de forma que se comprueben primero las restricciones gordas, esto es al generar sucesores.

Tecnica fordward chaining: (o fordward cheking, ha dicho ambos nombres no se cual es) quito todos los valores de las otras variables incompatibles con el padre, pero solo a los hijos de un cierto paso.



\section{Algoritmos}

Algoritmo de Satisfaccion de restricciones: 
\begin{itemize}
    \item AC3:
    \item 
\end{itemize}

Para el AC3 de la practica voy a ir reduciendo los dominios de las variables y voy a ir expandiendo  los dominios mas pequeños para buscar la solucion

\section{Ejecrcicios}
\subsection{Ejericio 1} Dada la red inicial de la figura, donde en cada nodo se representa el dominio de la
variable, obténgase la red arco consistente equivalente.
<Falta imagen>

Solucion: <Buscar una bien escrita>

Variables: X1,X2,X3
Restricciones: X1>X2>X3
Paso 1:
Dom(X1)={6,10}
Dom(X2)={5,9}
Dom(X3)={8,15}
Paso 2:
Dom(X1)={6,10}
Dom(X2)={9,9}
Dom(X3)={8,15}
Paso 3:
Dom(X1)={10,10}
Dom(X2)={9,9}
Dom(X3)={8,15}
Paso 4:
Dom(X1)={10,10}
Dom(X2)={9,9}
Dom(X3)={8,8}

\subsection{Ejercicio 2} El sudoku es un claro ejemplo de representaciópn del puzzle como un problema de
satisfacci´on de restricciones. Este problema se publicó en Nueva York en el año 1979
bajo el nombre de “Number Place” y se hizo popular en Japón con el nombre de sudoku
(“Sudji wa dokushin ni kagiru”: “los números deben ser sencillos” o “los números deben
aparecer una vez”).
\medskip
En el problema del sudoku (ver un ejemplo en la figura), hay que rellenar las casillas de
un tablero de 9 x 9 con números del 1 al 9, de forma que no se repita ningún número
en la misma fila, columna, o subcuadro de 3 x 3 que componen el sudoku. Modelar
este problema como un CSP.
<Falta imagen>

\subsection{Ejercicio 3} Cinco casas consecutivas tienen colores diferentes y son habitadas por hombres de
diferentes nacionalidades. Cada uno tiene un animal diferente, una bebida preferida y
fuma una marca determinada. Además, se sabe que:
\begin{itemize}
    \item El noruego vive en la primera casa
    \item La casa de al lado del noruego es azul
    \item El habitante de la tercera casa bebe leche
    \item El inglés vive en la casa roja
    \item El habitante de la casa verde bebe café
    \item El habitante de la casa amarilla fuma Kools
    \item La casa blanca se encuentra justo después de la verde 8. El español tiene un perro
    \item  El ucraniano bebe té
    \item El japonés fuma Cravens
    \item El fumador de Old Golds tiene un caracol
    \item El fumador de Gitanes bebe vino
    \item El vecino del fumador de Chesterfields tiene un reno
    \item El vecino del fumador de Kools tiene un caballo
\end{itemize}

\subsection{Ejercicio 4} Obtener la ordenaci´on de variables seg´un los diferentes m´etodos de ordenaci´on est´aticos
para cada uno de los CSP binarios siguientes:
<Falta imagen>

\subsection{Ejercicio 5} Seis amigos y sus respectivas esposas, se hospedan en el mismo hotel; y todos ellos
salen todos los d´ıas, asistiendo a reuniones de distinto volumen y composici´on. Para
asegurar la variedad en estas salidas, han acordado establecer las siguientes reglas: “Si
Antonio est´a con su mujer, es decir en la misma reuni´on que su mujer, y David con
la suya, y Luis con la se˜nora de Pedro, Enrique debe estar con la se˜nora de Ram´on.
Si Antonio est´a con su mujer y Pedro con la suya, y David con la se˜nora de Enrique,
Ram´on no debe estar con la se˜nora de Luis. Si Enrique y Ram´on y sus mujeres est´an
todos en la misma reuni´on, y Antonio no est´a con la se˜nora de David, Luis no debe
estar con la se˜nora de Pedro. Si Antonio est´a con su mujer y Ram´on con la suya, y
David no est´a con la se˜nora de Enrique, Luis debe estar con la se˜nora de Pedro. Si Luis
est´a con su mujer y Pedro con la suya y Enrique con la se˜nora de Ram´on, Antonio no
debe estar con la se˜nora de David. Si David y Enrique y sus mujeres est´an todos en la
misma reuni´on, y Luis no est´a con la se˜nora de Pedro, Ram´on debe estar con la se˜nora
de Luis”.
\medskip
Modelar dicho problema como un CSP y obtener las posibles combinaciones de reuniones para cada amigo. ¿Es posible que todos los d´ıas haya al menos un matrimonio
cuyos miembros no est´en juntos en la misma reuni´on? Este problema y el siguiente son
originales de Lewis Carroll, en su obra L´ogica Simb´olica.


\subsection{Ejercicio 6} Cinco amigos, Bernardo, Casimir, Luis, Carlos y Marcial, se encuentran cada d´ıa en el
restaurante. Tienen estas reglas, que observan cada vez que comen: “Bernardo toma
sal si y solamente si Casimir toma solo sal o solo mostaza. Bernardo toma mostaza si y
solamente si, o bien Luis no toma sal ni mostaza, o bien Marcial toma ambas. Casimir
toma sal si y solamente si, o Bernardo toma solamente uno de los dos condimentos, o
bien Marcial no toma ninguno de ellos. Toma mostaza si y solamente si Luis o Carlos
toman dos condimentos. Luis toma sal si y solamente si o bien Bernardo no toma
ning´un condimento, o bien Casimir toma ambos. Luis toma mostaza si y solamente
si Carlos o Marcial no toman ni sal ni mostaza. Carlos toma sal si y solamente si
Bernardo o Luis no toman ni sal ni mostaza. Toma mostaza si y solamente si Casimir
o Marcial no toman ni sal, ni mostaza. Marcial toma sal si y solamente si Bernardo o
Carlos toman los dos condimentos. Marcial toma mostaza si y solamente si Casimir o
Luis toman solo un condimento”.
\medskip
Modelar el CSP correspondiente. El problema consiste en descubrir si estas reglas son
compatibles y, en caso de que lo sean, cuales son las combinaciones posibles.

\subsection{Ejercicio 7} Representar mediante un CSP la siguiente informaci´on: Juana, Pepa y Paloma nacieron
y viven en ciudades diferentes (M´alaga, Madrid y Valencia). Adem´as, ninguna vive en
la ciudad donde naci´o. Juana es m´as alta que la que vive en Madrid. Paloma es cu˜nada
de la que vive en Valencia. La que vive en Madrid y la que naci´o en M´alaga tienen
nombres que comienzan por distinta letra. La que naci´o en M´alaga y la que vive ahora
en Valencia tienen nombres que comienzan por la misma letra.
\medskip
¿ D´onde naci´o y vive cada una?

\subsection{Ejercicio 8} Henry Dudeney (1847-1930) era un ingenioso inventor de problemas matem´aticos. Entre sus aportaciones se encuentra un puzzle con cierta complejidad de resoluci´on. Se
trata de un hept´agono donde en cada arista hay que colocar tres n´umeros: uno en cada
v´ertice y otro en el centro de la arista. Hay que colocar los n´umeros del 1 al 14 alrededor
de las aristas del hept´agono de manera que el n´umero de la arista y los dos v´ertices
extremos sumen lo mismo, para todas las aristas. Modelar el CSP correspondiente.
<Falta imagen>

\subsection{Ejercicio 9} Obtener, seg´un las heur´ısticas de ordenaci´on de variables, la mejor ordenaci´on para el
coloreado de los siguientes mapas:


<Falta imagen>


% ---------Bibliografía del Tema------------
\bigskip
\noindent\rule{\textwidth}{0.4pt} 

\section*{Bibliografía del Tema}
\addcontentsline{toc}{section}{Bibliografía del Tema}

\begin{itemize}
    \item
\end{itemize}
\newpage